\chapter{Baire空间和维数论}
\label{chapter:ch08}

本章我们介绍一类称为 Baire 空间的拓扑空间. 要陈述定义 Baire 空间的条件并不容易, 但它却在分析学和拓扑学中有广泛应用. 我们所学过的空间大多数都是 Baire 空间. 例如, 如果一个 Hausdorff 空间是紧致的或者是局部紧致的, 那么它便是一个 Baire 空间. 如果一个可度量化的空间 $X$ 是拓扑完备的, 即存在 $X$ 的一个度量, 关于这个度量 $X$ 是完备的, 那么它便是一个 Baire 空间.

因此, 从 $X$ 到 $\mathbb{R}^{n}$ 的所有连续函数构成的空间 $\mathcal{C}(X,\mathbb{R}^{n})$ 在一致度量下是完备的, 可见它关于一致拓扑是一个 Baire 空间. 这个结论有一系列有趣的应用.

其中一个应用就是我们将在第 49 节中给出的关于连续却处处不可微的实值函数存在性的证明.

另一个应用出现在拓扑学的一个分支维数论中. 第 50 节中我们将定义拓扑维数的概念, 它源自于 Lebesgue. 我们将证明一个经典的定理: 每一个拓扑维数为 $m$ 的紧致可度量化空间, 都可以嵌入到维数为 $N=2m+1$ 的欧氏空间 $\mathbb{R}^{N}$ 中. 由此可见, 每一个 $m$-流形都可以嵌入到 $\mathbb{R}^{2m+1}$ 中. 这就推广了第 36 节中证明的嵌入定理.

在本章中, 我们假设读者已熟悉完备度量空间 (第 43 节). 当学习维数论的时候, 将要用到第 36 节 (流形的嵌入) 以及一些线性代数的内容.

\subfile{Chapter-08/section-48}
\subfile{Chapter-08/section-49}
\subfile{Chapter-08/section-50}
\subfile{Chapter-08/section-extra}
